\documentclass[fleqn]{article}

\usepackage[a4paper,margin=1.5cm]{geometry}
\usepackage{amsmath,amssymb}
\usepackage{enumitem}
\usepackage{multicol}

\setlength{\mathindent}{0pt}
\setlist{nosep}

\begin{document}
\raggedright

\begin{multicols}{2}

\section*{Linear Algebra Cheat Sheet --- Compact \& Complete}

\textit{Convention: a matrix of size $r \times c$ has $r$ rows and $c$ columns.}

\section*{1. Fields}
A field $\mathbb{F}$ is a set with operations $(+,\cdot)$ satisfying: closure, associativity, commutativity, identities $(0,1)$, inverses, and distributivity.

\section*{2. Matrices}
\[
A_{m\times p} B_{p\times n} = (AB)_{m\times n},
\qquad
(AB)_{ij} = \sum_{k=1}^{p} a_{ik} b_{kj}
\]
\begin{itemize}
\item $AB = 0$ does not imply $A=0$ or $B=0$
\item If $A^T = A$ and $B^T = B$, then $AB$ is symmetric if and only if $AB = BA$
\item If $\varphi$ is an elementary row operation and $A \in \mathbb{F}^{m\times n}$, then
\[
\varphi(A) = \varphi(I_m)\,A
\]
\end{itemize}

\section*{3. Linear Systems}
\[
Ax=b, \qquad n \text{ variables}, \qquad r = \operatorname{rank}(A)
\]
\begin{itemize}
\item Number of pivots $= r$, free variables $= n-r$
\item System is solvable if and only if $b \in \operatorname{Col}(A)$
\item Homogeneous system: $Ax=0$ has solution space dimension $n-r$
\end{itemize}

\[
\{x : Ax=b\} = x_p + \ker(A)
\]

\section*{4. Rank}
\begin{itemize}
\item $\operatorname{rank}(A_{m\times n}) \le \min\{m,n\}$
\item $\operatorname{rank}(A^T)=\operatorname{rank}(A)$
\item $A \sim B \Rightarrow \operatorname{rank}(A)=\operatorname{rank}(B)$
\item $\operatorname{rank}(AB) \le \min\{\operatorname{rank}(A),\operatorname{rank}(B)\}$
\item $A,B \in \mathbb{F}^{n\times n},\; AB=0 \Rightarrow \operatorname{rank}(A)+\operatorname{rank}(B) \le n$
\end{itemize}

\section*{5. Vector Spaces \& Subspaces}
A subset $W \subseteq V$ is a subspace if and only if:
\begin{itemize}
\item $W$ is nonempty (equivalently $0 \in W$)
\item Closed under addition: if $u,v \in W$ then $u+v \in W$
\item Closed under scalar multiplication: if $\alpha \in \mathbb{F}$ and $u \in W$ then $\alpha u \in W$
\end{itemize}
Equivalently: $\alpha u + v \in W$ for all $u,v \in W$ and all $\alpha \in \mathbb{F}$.

\begin{itemize}
\item $U \cap W$ is a subspace of $V$
\item $U \cup W$ is a subspace of $V$ if and only if $U \subseteq W$ or $W \subseteq U$
\item $U+W = \{u+w : u \in U,\ w \in W\}$ is a subspace of $V$
\item Direct sum: $U \oplus W$ if and only if $U \cap W = \{0\}$
\item $\operatorname{span}(S)$ is the smallest subspace containing $S$
\end{itemize}

\section*{6. Bases \& Dimension}
If $\dim V = n$:
\begin{itemize}
\item Any linearly independent set has size at most $n$
\item Any spanning set has size at least $n$
\item A set $B$ with $|B|=n$ is a basis if and only if it is linearly independent if and only if it spans $V$
\item If $W \le V$, then $W=V$ if and only if $\dim W = \dim V$
\end{itemize}
\[
\dim(U+W) = \dim U + \dim W - \dim(U \cap W)
\]

\section*{7. Matrix Spaces}
\begin{itemize}
\item $\dim \operatorname{Row}(A) = \dim \operatorname{Col}(A) = \operatorname{rank}(A)$
\item $\dim \ker(A) = n - \operatorname{rank}(A)$ for $A \in \mathbb{F}^{m\times n}$
\item $\operatorname{Row}(AB) \subseteq \operatorname{Row}(B)$
\item $\operatorname{Col}(AB) \subseteq \operatorname{Col}(A)$
\end{itemize}

\section*{8. Invertible Matrices}

\begin{itemize}
\item If $P$ is invertible, then
\[
\operatorname{rank}(PA)=\operatorname{rank}(AP)=\operatorname{rank}(A)
\]
\item If $\alpha \neq 0$, then
\[
(\alpha A)^{-1}=\alpha^{-1}A^{-1}
\]
\item If $C$ is invertible and $B=CA$, then
\[
B \text{ invertible} \iff A \text{ invertible}
\]
\end{itemize}

\subsection*{Conditions Related to Invertibility}
Let $A \in \mathbb{F}^{n\times n}$. Each item states its relation to ``$A$ is invertible'':

\begin{enumerate}
\item $\;\Longrightarrow\; A^{-1}$ is unique
\item $\;\Longrightarrow\;$ $A$ has no zero row or column
\item $\iff \operatorname{rank}(A)=n$
\item $\iff \det(A)\neq 0$
\item $\iff A^T$ is invertible, with $(A^T)^{-1}=(A^{-1})^T$
\item $\iff A \sim I_n$
\item $\iff \operatorname{RREF}(A)=I_n$
\item $\iff Ax=0 \Rightarrow x=0$
\item $\iff$ For all $b\in\mathbb{F}^n$, the system $Ax=b$ has the unique solution $x=A^{-1}b$
\item $\iff \operatorname{Row}(A)=\mathbb{F}^n$
\item $\iff \operatorname{Col}(A)=\mathbb{F}^n$
\item $\iff$ The rows of $A$ span $\mathbb{F}^n$
\item $\iff$ The columns of $A$ span $\mathbb{F}^n$
\item $\iff$ The rows of $A$ are linearly independent
\item $\iff$ The columns of $A$ are linearly independent
\item $\iff$ The rows of $A$ form a basis of $\mathbb{F}^n$
\item $\iff$ The columns of $A$ form a basis of $\mathbb{F}^n$
\item $\iff \dim \operatorname{Row}(A)=n$
\item $\iff \dim \operatorname{Col}(A)=n$
\item $\iff A$ is a product of elementary matrices
\item $\iff \operatorname{adj}(A)$ is invertible
\item $\iff$ $A$ is invertible on one side (i.e.\ $AB=I$ or $BA=I$ for some $B$)
\item $\;\Longleftarrow\;$ $A$ is elementary (then $A^{-1}$ is also elementary)
\end{enumerate}

\section*{9. Determinants \& Adjugate}
\begin{itemize}
\item $\det(A^T)=\det(A)$
\item $\det(AB)=\det(A)\det(B)$
\item $\det(A^{-1}) = (\det A)^{-1}$
\item $\det(\alpha A)=\alpha^n \det(A)$ for $A \in \mathbb{F}^{n\times n}$
\end{itemize}
\[
A\,\operatorname{adj}(A)=\det(A)I,
\qquad
A^{-1} = \frac{1}{\det(A)}\operatorname{adj}(A)
\]
\[
\operatorname{rank}(\operatorname{adj}(A)) =
\begin{cases}
n, & \operatorname{rank}(A)=n, \\
1, & \operatorname{rank}(A)=n-1, \\
0, & \operatorname{rank}(A) < n-1
\end{cases}
\]

\section*{10. Linear Maps}
Let $T:V \to U$ be linear.
\begin{itemize}
\item Kernel: $\ker T = \{v \in V : T(v)=0\}$
\item Image: $\operatorname{Im} T = \{T(v) : v \in V\}$
\item $T$ is injective if and only if $\ker T = \{0\}$
\item $T$ is surjective if and only if $\dim \operatorname{Im}T = \dim U$
\end{itemize}
\[
\dim V = \dim \ker T + \dim \operatorname{Im}T
\]

\section*{11. Operators \& Isomorphisms}
\begin{itemize}
\item A linear map is an isomorphism if and only if it is injective and surjective
\item $V \cong U$ if and only if $\dim V = \dim U$
\item For an operator $T:V \to V$:
\[
\ker T \subseteq \ker T^2,
\qquad
\operatorname{Im}T \supseteq \operatorname{Im}T^2
\]
\item In finite dimensions, equality implies stabilization
\end{itemize}

\section*{12. Matrix Representation}
Let $B=\{v_1,\dots,v_n\}$ and $S=\{u_1,\dots,u_m\}$.
\[
[T]_B^S = \big([T(v_1)]_S, \dots, [T(v_n)]_S\big)
\]
\begin{itemize}
\item $[T]_B^S [v]_B = [T(v)]_S$
\item $[S\circ T]_A^C = [S]_B^C [T]_A^B$
\item $T$ is invertible if and only if $[T]_B^S$ is invertible
\end{itemize}

\section*{13. Dual Spaces, Functionals, and Annihilators}

\textbf{Concept.}
Duality studies vectors indirectly via linear measurements. Instead of examining vectors themselves, we examine how all linear functionals act on them.

\textbf{Linear functionals.}
A linear functional is a linear map $f:V \to \mathbb{F}$. It assigns a scalar to each vector and represents a linear measurement.

\textbf{Dual space.}
\[
V^* = \{ f: V \to \mathbb{F} \mid f \text{ is linear} \}
\]
$V^*$ is a vector space. If $\dim V = n < \infty$, then $\dim V^* = n$.

\textbf{Dual basis.}
If $B = \{v_1,\dots,v_n\}$ is a basis of $V$, the dual basis
$B^* = \{f_1,\dots,f_n\}$ satisfies $f_i(v_j)=\delta_{ij}$.
Each $f_i$ extracts the $i$-th coordinate.

\textbf{Second dual space.}
The second dual is $V^{**} = (V^*)^*$.
There is a canonical map
\[
\Phi:V \to V^{**},
\qquad
\Phi(v)(f)=f(v)
\]
In finite dimensions, this map is an isomorphism.

\textbf{Kernel (MERHAV HA EFES).}
For $T:V \to U$,
\[
\ker T = \{v \in V : T(v)=0\}
\]
This is a subspace of $V$.

\textbf{Annihilator (MEAPES).}
For $W \le V$,
\[
W^{\circ} = \{f \in V^* : f(w)=0 \text{ for all } w \in W\}
\]
This is a subspace of $V^*$.

\textbf{Difference between MERHAV HA EFES and MEAPES.}
\begin{itemize}
\item MERHAV HA EFES: vectors in $V$ mapped to zero by a linear map
\item MEAPES: functionals in $V^*$ that vanish on a subspace
\end{itemize}

\textbf{Dimension relation.}
If $W \le V$ and $\dim V = n < \infty$, then
\[
\dim W + \dim W^{\circ} = n
\]

\textbf{Dual (transpose) map.}
For $T:V \to U$ define
\[
T^*:U^* \to V^*,
\qquad
T^*(f)=f \circ T
\]
Then
\[
\ker T^* = (\operatorname{Im}T)^{\circ},
\qquad
\operatorname{Im}T^* = (\ker T)^{\circ}
\]
In matrix form,
\[
[T^*] = [T]^T
\]

\section*{14. Vandermonde Matrix}

Given scalars $x_1,\dots,x_n \in \mathbb{F}$, the \textbf{Vandermonde matrix} is
\[
V(x_1,\dots,x_n)=
\begin{pmatrix}
1 & x_1 & x_1^2 & \cdots & x_1^{\,n-1} \\
1 & x_2 & x_2^2 & \cdots & x_2^{\,n-1} \\
\vdots & \vdots & \vdots & \ddots & \vdots \\
1 & x_n & x_n^2 & \cdots & x_n^{\,n-1}
\end{pmatrix}
\]

\textbf{Determinant.}
\[
\det V(x_1,\dots,x_n)
=
\prod_{1 \le i < j \le n} (x_j - x_i)
\]

\textbf{Invertibility.}
\[
V(x_1,\dots,x_n) \text{ is invertible }
\iff
x_1,\dots,x_n \text{ are pairwise distinct}
\]

\textbf{Consequences.}
\begin{itemize}
\item If $x_i = x_j$ for some $i \neq j$, then two rows are equal $\Rightarrow \det V = 0$
\item If all $x_i$ are distinct, then $\operatorname{rank}(V)=n$
\item Appears naturally in polynomial interpolation
\end{itemize}

\section*{15. Permutations}

A \textbf{permutation} of $\{1,\dots,n\}$ is a bijection
\[
\sigma:\{1,\dots,n\}\to\{1,\dots,n\}.
\]
The set of all permutations is denoted $S_n$.

\textbf{Inversions.}
An inversion of $\sigma$ is a pair $(i,j)$ such that
\[
i<j \quad \text{and} \quad \sigma(i)>\sigma(j).
\]
Let $\operatorname{inv}(\sigma)$ be the number of inversions of $\sigma$, sometimes denoted $S(\sigma)$.

\textbf{Sign of a permutation.}
\[
\operatorname{sgn}(\sigma)=(-1)^{\operatorname{inv}(\sigma)}.
\]
\begin{itemize}
\item $\sigma$ is \emph{even} if $\operatorname{sgn}(\sigma)=1$
\item $\sigma$ is \emph{odd} if $\operatorname{sgn}(\sigma)=-1$
\end{itemize}

\textbf{Properties.}
\begin{itemize}
\item $\operatorname{sgn}(\sigma\tau)=\operatorname{sgn}(\sigma)\operatorname{sgn}(\tau)$
\item A transposition changes the sign
\item $|S_n| = n!$
\end{itemize}

\textbf{Determinant (Leibniz formula).}
For $A=(a_{ij}) \in \mathbb{F}^{n\times n}$,
\[
\det(A)=
\sum_{\sigma\in S_n}
\operatorname{sgn}(\sigma)
\prod_{i=1}^n a_{i,\sigma(i)}.
\]

\textbf{Consequences.}
\begin{itemize}
\item If two rows (or columns) are equal, every term cancels $\Rightarrow \det(A)=0$
\item Swapping two rows multiplies $\det(A)$ by $-1$
\item Linearity of determinant follows from the formula
\end{itemize}

\section*{16. Block Matrices}

\textbf{Definition.}
A \emph{block matrix} is a matrix partitioned into submatrices:
\[
A =
\begin{pmatrix}
A_{11} & A_{12} \\
A_{21} & A_{22}
\end{pmatrix},
\]
where the blocks have compatible sizes.

\textbf{Block operations.}
If dimensions agree, block matrices behave like ordinary matrices:
\begin{itemize}
\item Addition and scalar multiplication are blockwise
\item Matrix multiplication follows block multiplication rules
\end{itemize}
\[
\begin{pmatrix}
A & B\\
C & D
\end{pmatrix}
\begin{pmatrix}
E & F\\
G & H
\end{pmatrix}
=
\begin{pmatrix}
AE+BG & AF+BH\\
CE+DG & CF+DH
\end{pmatrix}
\]

\textbf{Block diagonal matrices.}

If
\[
A=
\begin{pmatrix}
A_1 & 0\\
0 & A_2
\end{pmatrix},
\]
\[
\begin{aligned}
\det(A) &= \det(A_1)\det(A_2), \\
A \text{ invertible } &\iff A_1, A_2 \text{ invertible}.
\end{aligned}
\]


\textbf{Block triangular matrices.}
If
\[
A=
\begin{pmatrix}
A_{11} & A_{12}\\
0 & A_{22}
\end{pmatrix},
\]
then
\[
\det(A)=\det(A_{11})\det(A_{22}),
\qquad
\operatorname{rank}(A)=\operatorname{rank}(A_{11})+\operatorname{rank}(A_{22}).
\]

\textbf{Schur complement.}
Let
\[
A=
\begin{pmatrix}
A_{11} & A_{12}\\
A_{21} & A_{22}
\end{pmatrix},
\quad A_{11} \text{ invertible}.
\]
The \emph{Schur complement} of $A_{11}$ in $A$ is
\[
S = A_{22} - A_{21}A_{11}^{-1}A_{12}.
\]

\begin{itemize}
\item $A$ invertible $\iff A_{11}$ and $S$ invertible
\item $\det(A)=\det(A_{11})\det(S)$
\end{itemize}

\textbf{Block inverse formula.}
If $A_{11}$ and $S$ are invertible, then
\[
A^{-1} =
\begin{pmatrix}
A_{11}^{-1} + A_{11}^{-1}A_{12}S^{-1}A_{21}A_{11}^{-1}
&
- A_{11}^{-1}A_{12}S^{-1}
\\[6pt]
- S^{-1}A_{21}A_{11}^{-1}
&
S^{-1}
\end{pmatrix}.
\]

\textbf{Solving block systems.}
For
\[
\begin{pmatrix}
A_{11} & A_{12}\\
A_{21} & A_{22}
\end{pmatrix}
\begin{pmatrix}
x\\
y
\end{pmatrix}
=
\begin{pmatrix}
b\\
c
\end{pmatrix},
\]
solve:
\begin{enumerate}
\item $A_{11}x = b - A_{12}y$
\item $(A_{22}-A_{21}A_{11}^{-1}A_{12})y = c - A_{21}A_{11}^{-1}b$
\end{enumerate}

\textbf{Implementation tips.}
\begin{itemize}
\item Use block structure to reduce computation
\item Common in Gaussian elimination, LU decomposition, and numerical methods
\item Essential for large sparse systems and matrix factorizations
\end{itemize}


\end{multicols}
\end{document}

